% Part: lambda-calculus
% Chapter: lambda-definability
% Section: introduction

\documentclass[../../../include/open-logic-section]{subfiles}

\begin{document}

\olfileid[pt]{lam}{ldf}{int}
\olsection{Introdução}

À primeira vista, o cálculo lambda é apenas um cálculo muito abstrato de
expressões que representam funções e aplicações delas a outras. Nada na
sintaxe do cálculo lambda sugere que sejam funções de tipos específicos de
objetos; em particular, a sintaxe não menciona números naturais. Suas
operações básicas---aplicação e abstrações lambda---são operações aplicáveis
a qualquer função, não apenas a funções nos números naturais.

Ainda assim, com alguma engenhosidade, é possível definir funções
aritméticas, isto é, funções nos números naturais, no cálculo lambda. Para
isso, definimos, para cada número natural~$n \in \Nat$, um termo-$\lambd$
especial~$\num n$, o \emph{numeral de Church} para~$n$. (Os numerais de
Church recebem esse nome em homenagem a Alonzo Church.)

\begin{defn}
  Se $n \in \Nat$, o \emph{numeral de Church} correspondente $\num{n}$
  representa~$n$:
  \[
    \num{n} \ident \lambd[fx][f^n(x)]
  \]
  Aqui, $f^n(x)$ representa o resultado de aplicar $f$ a~$x$ $n$ vezes. Por
  exemplo, $\num{0}$ é $\lambd[fx][x]$, e $\num{3}$ é
  $\lambd[fx][f(f(f\,x))]$.
\end{defn}
  
O numeral de Church $\num n$ é codificado como um termo lambda que
representa uma função que aceita dois argumentos $f$ e $x$ e devolve
$f^n(x)$. Os numerais de Church estão evidentemente em forma normal.

Uma representação dos números naturais no cálculo lambda só é útil, é
claro, se pudermos calcular com eles. Calcular com numerais de Church no
cálculo lambda significa aplicar um termo-$\lambd$~$F$ a tal numeral de
Church e reduzir o termo combinado~$F\, \num n$ a uma forma normal. Se ele
sempre se reduz a uma forma normal, e essa forma normal é sempre um numeral
de Church~$\num m$, podemos considerar que a saída do cálculo é o número~$m$.
Podemos então considerar que~$F$ define uma função
$f\colon \Nat \to \Nat$, a saber, a função tal que $f(n) = m$ sse
$F\, \num n \red \num m$. Pela propriedade de Church--Rosser, formas
normais são únicas quando existem. Assim, se $F\, \num n \red \num m$, não
pode haver outro termo em forma normal, em particular outro numeral de
Church, ao qual $F \, \num n$ se reduza.

Reciprocamente, dada uma função $f\colon \Nat \to \Nat$, podemos perguntar
se existe um termo $F$ que defina~$f$ dessa maneira. Nesse caso, dizemos que
$F$ \emph{!!{lambda define}s}~$f$ e que $f$ é !!{lambda definable}. Podemos
generalizar isso para funções de vários argumentos e funções parciais.

\begin{defn}
Suponha $f\colon \Nat^k \to \Nat$. Dizemos que um termo lambda~$F$
\emph{!!{lambda define}s} $f$ se, para todos
$n_0$, \dots,~$n_{k-1}$,
\[
F \, \num{n_0} \, \num{n_1} \dots \num{n_{k-1}} \red \num{f(n_0, n_1, \dots,
  n_{k-1})}
\]
se $f(n_0, \dots, n_{k-1})$ está definida, e $F \, \num{n_0} \,
\num{n_1} \dots \num{n_{k-1}}$ não tem forma normal caso contrário.
\end{defn}

Um exemplo muito simples são as funções constantes. O termo $C_k \ident
\lambd[x][\num{k}]$ !!{lambda define}s a função $c_k\colon \Nat \to
\Nat$ tal que $c(n) = k$. Pois $C_k \, \num n \ident
(\lambd[x][\num{k}])\num n \redone \num{k}$ para qualquer~$n$. A função
identidade é !!{lambda defined} por $\lambd[x][x]$. Funções mais complexas
são, naturalmente, mais difíceis de definir e muitas vezes exigem bastante
engenhosidade. Portanto, talvez seja surpreendente que toda função
computável seja !!{lambda definable}. A recíproca também é verdadeira: se
uma função é !!{lambda definable}, ela é computável.

\end{document}

